\scanpage{217}
\mspage{756}

this group scheme $\underline{\Pi}^{\circ}$ over $S_0=\operatorname{Spec}\mathbb{Z}$, as an étale sheaf on the ``big étale topos'' of $S_0$, is generated by \deleted{the}\inserted{a finite number of} subgroups $k_i'(\underline{\mathrm{G}_a})$,
\deleted{and by the discrete group $\Pi_0$.}
\inserted{(with possible repetition of the same conjugacy class), or again \deleted{\unsure{rather}}\inserted{that} $\underline{\Pi}$ is}
\deleted{(And there \unsure{arises} \unsure{infinitely|the interesting} \unsure{questi-}}
\inserted{generated by \unsure{the}}
the $k_i'(\mathrm{G}_a)$ (\unsure{their various conjugates|one of each} say) and the discrete group $(\Pi_0)_{S_0}$.
(And there arises the interesting question of making the relations explicit\ldots{} is it enough to write the relations of $\Pi_0$, plus \unsure{the commutativity} of
\[
(57)\qquad
\diagram{\begin{tikzcd}[ampersand replacement=\&,column sep=large,row sep=large]
T_0 \arrow[r,no head,"{k_i'}"] \arrow[d,"{\mathrm{can}}"] \& \Pi_0 \arrow[d,"{\mathrm{can}}"] \\
\underline{T_0} \arrow[r,"{\underline{k_i'}}"] \& \underline{\Pi}
\end{tikzcd}}
\qquad ?\quad\Big)
\]

Other conditions, concerning the group schemes
\[
(58)\qquad
\underline{\gamma}_{\underline{\Pi}}=\underline{\Pi}/\underline{\Pi}^{\circ},\quad
\underline{\gamma}_{\underline{\mathcal{G}}}=\underline{\mathcal{G}}/\underline{\mathcal{G}}^{\circ},\quad
\underline{\gamma}_{\underline{M}}=\underline{M}/\underline{M}^{\circ}
\]
-- we would like these to be \uline{constant} group schemes, \unsure{corresponding} to the \deleted{\unsure{schemes}} finite discrete groups which describe the corresponding quotients of the groups over $\mathbb{Q}$, which discrete groups are quotients
\margin{216}


\scanpage{218}

\[
(59)\qquad
\left\{\begin{array}{l}
\gamma_{\Pi}\ (\text{resp. }\gamma_{\mathcal{G}},\ \text{resp. }\gamma_M)\quad\text{quotients of }\Pi,\mathcal{G},M\\
\underline{\gamma}_{\underline{\Pi}}\simeq(\gamma_{\Pi})_{S_0},\quad
\underline{\gamma}_{\underline{\mathcal{G}}}\simeq(\gamma_{\mathcal{G}})_{S_0},\quad
\underline{\gamma}_{\underline{M}}\simeq(\gamma_M)_{S_0}
\end{array}\right.
\]
which fit into a sequence \deleted{exact}\inserted{of homomorphisms}
\[
(60)\qquad
\diagram{\begin{tikzcd}[ampersand replacement=\&,column sep=large]
\gamma_{\Pi} \arrow[r] \& \gamma_{\mathcal{G}} \arrow[r] \& \gamma_M
\end{tikzcd}}
\]
and into exact sequences
\[
(61)\qquad
\left\{\diagram{\begin{tikzcd}[ampersand replacement=\&,column sep=large,row sep=small]
\gamma_{\Pi} \arrow[r,no head] \& \gamma_{\mathcal{G}} \arrow[r] \& \gamma_{\unsure{\mathcal{T}}} \arrow[r] \& 1\\
\gamma_{\Pi} \arrow[r] \& \gamma_M \arrow[r] \& \gamma_N \arrow[r] \& 1\\
\gamma_{\mathcal{G}} \arrow[r] \& \gamma_M \arrow[r,no head] \& \gamma_{\deleted{\unsure{\mathcal{T}}}\Gamma} \arrow[r] \& 1\\
\gamma_{\unsure{\mathcal{T}}} \arrow[r,no head] \& \gamma_N \arrow[r] \& \gamma_{\Gamma} \arrow[r] \& 1
\end{tikzcd}}\right.\quad .
\]

This hypothesis on $\underline{\Pi},\underline{\mathcal{G}},\underline{M}$ makes it possible to define groups $\underline{\Pi}^{\Delta}(\widehat{\mathbb{Z}})^{\unsure{\Delta}},\underline{\mathcal{G}}^{\Delta}(\widehat{\mathbb{Z}})^{\unsure{\Delta}},\underline{M}^{\Delta}(\widehat{\mathbb{Z}})$, and we shall necessarily have the first two factorizations of
\[
(62)\qquad
\left\{\diagram{\begin{tikzcd}[ampersand replacement=\&,column sep=large,row sep=large]
\Pi \arrow[r] \& \underline{\Pi}^{\Delta}(\deleted{\unsure{\widehat{\mathbb{Q}}}}\widehat{\mathbb{Z}})\\
\mathcal{G} \arrow[r] \& \unsure{\underline{\mathcal{G}}}^{\Delta}(\widehat{\mathbb{Z}})\\
M \arrow[r] \& \underline{M}^{\Delta}(\widehat{\mathbb{Z}})
\end{tikzcd}}\right.
\]
(for
\[
\diagram{\begin{tikzcd}[ampersand replacement=\&,column sep=large,row sep=small]
\Pi_0 \arrow[r,no head] \& \underline{\Pi}(\mathbb{Z})\subset\underline{\Pi}^{\Delta}(\widehat{\mathbb{Z}})\\
\mathcal{G}_0 \arrow[r,no head] \& \underline{\mathcal{G}}(\mathbb{Z})\subset\underline{\mathcal{G}}^{\Delta}(\widehat{\mathbb{Z}})
\end{tikzcd}}
\qquad).
\]
As for the third factorization, it should be included in the starting conditions for a Lie homomorphism of a Teichmüller - Galois structure - namely that $\underline{\unsure{M}}_{\mathbb{Q}}/\underline{\unsure{M}}_{\mathbb{Q}}^{\circ}$ is constant, which makes it possible
\margin{\deleted{\unsure{smooth}}}
\margin{217}


\scanpage{219}
\mspage{757}

then to define $\underline{\unsure{M}}_{\mathbb{Q}}^{\Delta}(\unsure{A|\mathbb{A}}_{\mathbb{Q}})$, and we require that $M\relbar\underline{M}_{\mathbb{Q}}(\unsure{A|\mathbb{A}}_{\mathbb{Q}})$ factor through
\[
(63)\qquad
\diagram{\begin{tikzcd}[ampersand replacement=\&,column sep=large]
M \arrow[r] \& \underline{M}_{\mathbb{Q}}^{\Delta}(\unsure{A|\mathbb{A}}_{\mathbb{Q}})
\end{tikzcd}}\quad .
\]

If moreover we assume that the \unsure{images} of the three homomorphisms (62) are \uline{open} \unsure{subgroups}, (a condition which considerably strengthens the \unsure{index} condition \unsure{given} $4^\circ$) (p. 749), provided we combine it with the surjectivity of the composites
\[
(64)\qquad
\left\{\diagram{\begin{tikzcd}[ampersand replacement=\&,column sep=large,row sep=large]
\Pi \arrow[r,no head] \& \underline{\Pi}^{\Delta}(\widehat{\mathbb{Z}}) \arrow[r] \& \Gamma\underline{\gamma}_{\underline{\Pi}}\\
\mathcal{G} \arrow[r] \& \underline{\mathcal{G}}^{\Delta}(\widehat{\mathbb{Z}}) \arrow[r,no head] \& \Gamma\underline{\gamma}_{\underline{\mathcal{G}}}\\
M \arrow[r,no head] \& \underline{M}^{\Delta}(\widehat{\mathbb{Z}}) \arrow[r,no head] \& \Gamma\underline{\gamma}_{\underline{M}}
\end{tikzcd}}\right.\qquad),
\]
we shall be particularly happy. We shall say that we have a \uline{Lie \deleted{\unsure{homomorphism}}\inserted{approximation}} (over $\mathbb{Z}$, hence ``integral'') \deleted{perfect} non\unsure{degenerate}, if all the additional conditions considered previously are satisfied\ldots
\margin{218}

\scanpage{220}
\uline{Remark} In the preceding presentation, the group $\Pi_0$ \inserted{(or $\unsure{\Pi}$)} itself seems to recede into the background, in favor of $\mathcal{G}_0$ \inserted{(or $\mathcal{G}$)} and its $\mathrm{LT}$-structure. It seems to me indeed that the importance of the role of $\Pi_0$ fades behind that of the \inserted{\unsure{sub-}}group \deleted{\unsure{created|side}} of finite index $\Pi_0'$ (if one exists) which would be \deleted{generated by the $k_i'(T_0)\subset\Pi_0$ (and their} a group with loops and \inserted{of which} the $k_i'(T_0)$ \deleted{\unsure{and}} would be \deleted{conjugates, i.e. $\unsure{\Pi_0''}$, and its \unsure{closure}} exactly the loop subgroups. It is not \inserted{too} clear, even if such a $\Pi_0'$ exists, that it is necessaril\textsuperscript{y} unique (\unsure{let us take|think of} the case for example where $I=\varnothing$ !). The case which had occupied me until now, $\Pi_0=\mathcal{G}_0=\mathcal{G}_{0,3}^{+}\unsure{=|\text{--}}\unsure{\mathcal{G}_{1,1}^{+}|\mathcal{T}_{1,1}^{+}}$, has been a little ``misleading'', for \unsure{we then took} $\Pi_0:=\mathcal{G}_0$ \deleted{\unsure{this}} \unsure{precisely}. If we had thought \deleted{\unsure{As we}} \inserted{\unsure{in desire}} \unsure{right} at the beginning of \unsure{my work}, where \inserted{\unsure{by taking only}} $\Pi_0=\Pi_{0,3}$ as the starting group, we would have seen (as was moreover actually the case) the role of $\Pi_{0,3}$ dwindle in favor of the smaller subgroups (type $\Pi^{\unsure{\sharp}}$, giving rise to the $\mathcal{G}^{\unsure{\leftrightarrow}}$ of \S\,49 $\overline{\underline{\mathrm{VI}}}$\ldots)
\margin{219}

\scanpage{221}
\mspage{758}
$4^\circ)$ \uline{Lie} \deleted{\uline{homomorphism}}\inserted{approximation} \inserted{of $(\Pi_0,\mathcal{G}_0,M)$} \uline{induced by an approximation of} $\mathcal{G}_0$.

Let us start with a system $(\mathcal{G}_0\supset\Pi_0)$ of a \inserted{discrete} $\mathrm{LT}$-group $\mathcal{G}_0$, and of an invariant \inserted{$L$-}subgroup $\Pi_0$ of $\mathcal{G}_0$. Suppose \unsure{given} group \unsure{scheme} $\underline{\mathcal{G}}_{\mathbb Q}$ over $\mathbb Q$, and a homomorphism
\margin{We assume that the homomorphism $\unsure{\mathcal{G}_0}\longrightarrow\mathbb Z^*=\{\pm1\}$ defined by the action of $\mathcal{G}_0$ on the $L$-group $\Pi_0$ (\unsure{multipliers}) is trivial\ldots}
\diagram{\begin{tikzcd}[ampersand replacement=\&,column sep=large]
(65)\quad\varphi:\ \mathcal{G}_0 \arrow[r] \& \underline{\mathcal{G}}(\mathbb Q)
\end{tikzcd}}
having the property that for $\forall$ homomorphism
\diagram{\begin{tikzcd}[ampersand replacement=\&,column sep=large]
(66)\quad k_i:\ T_0 \arrow[r] \& \mathcal{G}_0
\end{tikzcd}}
which \unsure{is} part of the $\mathrm{LT}$-structure, the composite homomorphism
\diagram{\begin{tikzcd}[ampersand replacement=\&,column sep=large]
T_0 \arrow[r] \& \mathcal{G}_0 \arrow[r,no head] \& \underline{\mathcal{G}}(\mathbb Q)
\end{tikzcd}}
is quasi-unipotent. We assume moreover that $\underline{\mathcal{G}}_{\mathbb Q}$ is \uline{generated} by the set of \deleted{\unsure{its}} $\mathbb Q$-rational points \unsure{defined} by the elements of $\mathcal{G}_0$ via (65).

We consider $\underline{\mathcal{G}}_{\mathbb Q}$ as equipped with the homomorphism (65) and with the family of outer homomorphisms
\diagram{\begin{tikzcd}[ampersand replacement=\&,column sep=large]
(67)\quad\underline{k}_i:\ \underline{T_0}_{\mathbb Q} \arrow[r] \& \underline{\mathcal{G}}_{\mathbb Q}
\end{tikzcd}}
\margin{220}

\scanpage{222}
giving rise to compatibilities
\diagram{\begin{tikzcd}[ampersand replacement=\&,column sep=huge,row sep=large]
T_0 \arrow[r,"k_i"] \arrow[d,"\mathrm{can}"',"{\deleted{\unsure{\varphi}}}"] \& \mathcal{G}_0 \arrow[d,"{(65)}"] \\
(68)\quad\underline{T_0}_{\mathbb Q}(\mathbb Q) \arrow[r,"\underline{k}_i(\mathbb Q)"] \arrow[r,shorten >=12pt] \& \underline{\mathcal{G}}(\mathbb Q).
\end{tikzcd}}

I then consider the group scheme over $\mathbb Q$ of \inserted{\deleted{pairs} pairs $(u,\mu)$ of an} \unsure{automorphisms|automorphism} of $\underline{\mathcal{G}}_{\mathbb Q}$, and of a $\mu\in\Gamma G_{\mathrm m}$, \unsure{such|such} that for $\forall\ i\in I$, \deleted{we have} there exists locally a section $f_i$ of $\underline{\Pi}_{\mathbb Q}$ (generated by the image of $\Pi_0$ in $\underline{\mathcal{G}}(\mathbb Q)$), \unsure{such|such} that we have
\[
(69)\qquad u\bigl(\underline{k}'_i(\lambda)\bigr)=\operatorname{int}(f_i)^{-1}(\mu\lambda).
\]
I denote by $\underline{M}_{\mathbb Q}^{\varphi}$ this group scheme over $\mathbb Q$, augmented to $G_{\mathrm m\mathbb Q}$ by $(u,\mu)\longmapsto\mu$, by a homomorphism which is trivial on $\operatorname{Im}(\underline{\mathcal{G}}_{\mathbb Q}\to\underline{M}_{\mathbb Q}^{\varphi})$, hence factors through a homomorphism
\diagram{\begin{tikzcd}[ampersand replacement=\&,column sep=large]
(70)\quad\underline{\Gamma}_{\mathbb Q}^{\varphi}\overset{\mathrm{déf}}{=}\underline{M}_{\mathbb Q}^{\varphi}/\operatorname{Im}\underline{\mathcal{G}}_{\mathbb Q} \arrow[r,"\chi_{\underline{\Gamma}}"] \& G_{\mathrm m\mathbb Q}.
\end{tikzcd}}
If now $(\mathcal{G}_0,\Pi_0)$ is part of a Galois-Teichmüller system
\margin{221}